% GENERATED FILE. Edit canonical lessons, then run npm run generate:books.
% canonical-source-hash: fnv1a64:19cdaf79d6c0c18d
% canonical-lessons: SA-C276-sah-gatavan, SA-C276-sah-pathati-sma, SA-C276-ghatate, SA-S226-letter-gha, SA-R276-what-happened

\chapter{What Happened, and What Was Going On}
\label{ch:sa-what-happened-and-what-was-going-on}
\hlchaptermodality{276}

\begin{chapteropening}
\textbf{By the end of this chapter:} \emph{I can say what someone did, with the past that follows the doer (saḥ gatavān, sā gatavatī), and what was going on at the time (saḥ paṭhati sma; yadā … tadā …).}

Complete SA-R276-what-happened: say what happened and what was going on when it happened.
\end{chapteropening}

\section[\texorpdfstring{\sk{सः} \sk{गतवान्} (saḥ gatavān)}{सः गतवान् (saḥ gatavān)}]{\sk{सः} \sk{गतवान्} (saḥ gatavān) — the past that follows the doer}

\label{lesson:SA-C276-sah-gatavan}



\begin{quote}
Before the new one: say \textquotedblleft{}the boy went\textquotedblright{}, with the past you know. (\emph{bālaḥ agacchat}.) Then say \textquotedblleft{}yesterday\textquotedblright{}. (\emph{hyaḥ}.) This lesson gives a second past, the one the stories in this book are told in.
\end{quote}

\subsection*{You'll want to know: \sk{गतवान्}, \sk{गतवती}}
\begin{quote}
\textbf{\sk{सः} \sk{गतवान्।}}
\end{quote}

\begin{quote}
\emph{saḥ gatavān.} — \textquotedblleft{}He went.\textquotedblright{}
\end{quote}

\begin{quote}
\textbf{\sk{सा} \sk{गतवती।}}
\end{quote}

\begin{quote}
\emph{sā gatavatī.} — \textquotedblleft{}She went.\textquotedblright{}
\end{quote}

\textbf{\sk{गतवान्}} (\emph{gatavān}), \textquotedblleft{}went\textquotedblright{}, is a describing word made from the verb: \emph{-vān} for a man, \emph{-vatī} for a woman.

\begin{quote}
\textbf{\sk{सः} \sk{ह्यः} \sk{गृहं} \sk{गतवान्।}}
\end{quote}

\begin{quote}
\emph{saḥ hyaḥ gṛhaṃ gatavān.} — \textquotedblleft{}He went home yesterday.\textquotedblright{}
\end{quote}

\begin{quote}
\textbf{\sk{मीरा} \sk{ह्यः} \sk{आपणं} \sk{गतवती।}}
\end{quote}

\begin{quote}
\emph{mīrā hyaḥ āpaṇaṃ gatavatī.} — \textquotedblleft{}Mira went to the shop yesterday.\textquotedblright{}
\end{quote}

For several people, \textbf{\sk{गतवन्तः}} (\emph{gatavantaḥ}); for two, \textbf{\sk{गतवन्तौ}} (\emph{gatavantau}):

\begin{quote}
\textbf{\sk{ते} \sk{ह्यः} \sk{ग्रामं} \sk{गतवन्तः।}}
\end{quote}

\begin{quote}
\emph{te hyaḥ grāmaṃ gatavantaḥ.} — \textquotedblleft{}They went to the village yesterday.\textquotedblright{}
\end{quote}

It needs no \textbf{\sk{अस्मि}}: a man says \emph{ahaṃ gatavān}, \textquotedblleft{}I went\textquotedblright{}, a woman \emph{ahaṃ gatavatī}.

\subsection*{You'll want to know: the stem under it}
The part before \emph{-vat-} is the verb's \emph{-ta} form, which describes on its own: \textbf{\sk{गतः}} (\emph{gataḥ}), \textquotedblleft{}gone\textquotedblright{}. Each verb has its own, given with it in the dictionary:

\noindent
\begin{tabularx}{\linewidth}{@{}>{\raggedright\arraybackslash}X>{\raggedright\arraybackslash}X@{}}
\toprule
\textbf{the verb} & \textbf{the past (a man)} \\
\midrule
\textquotedblleft{}does\textquotedblright{} & \textbf{\sk{कृतवान्}} (\emph{kṛtavān}) \\
\textbf{\sk{पठति}} & \textbf{\sk{पठितवान्}} (\emph{paṭhitavān}) \\
\textbf{\sk{खादति}} & \textbf{\sk{खादितवान्}} (\emph{khāditavān}) \\
\textbf{\sk{लिखति}} & \textbf{\sk{लिखितवान्}} (\emph{likhitavān}) \\
\textbf{\sk{पश्यति}} & \textbf{\sk{दृष्टवान्}} (\emph{dṛṣṭavān}) \\
\textbf{\sk{आगच्छति}} & \textbf{\sk{आगतवान्}} (\emph{āgatavān}) \\
\bottomrule
\end{tabularx}

The thing done stays as it was:

\begin{quote}
\textbf{\sk{सः} \sk{अन्नं} \sk{खादितवान्।}}
\end{quote}

\begin{quote}
\emph{saḥ annaṃ khāditavān.} — \textquotedblleft{}He ate the food.\textquotedblright{}
\end{quote}

\begin{quote}
\textbf{\sk{मीरा} \sk{पुस्तकं} \sk{पठितवती।}}
\end{quote}

\begin{quote}
\emph{mīrā pustakaṃ paṭhitavatī.} — \textquotedblleft{}Mira read the book.\textquotedblright{}
\end{quote}

\textquotedblleft{}Said\textquotedblright{} comes from the other \textquotedblleft{}say\textquotedblright{} root, √vac: \textbf{\sk{उक्तवान्}} (\emph{uktavān}; Monier-Williams gives \emph{ukta} there).

\begin{quote}
\textbf{\sk{सा} \sk{किम्} \sk{उक्तवती}?}
\end{quote}

\begin{quote}
\emph{sā kim uktavatī?} — \textquotedblleft{}What did she say?\textquotedblright{}
\end{quote}

The past you know still stands: \emph{agacchat} and \emph{gatavān} both say \textquotedblleft{}went\textquotedblright{}, and both are correct. This book tells its stories with \emph{gatavān}, one shape for every verb (Pāṇini 1.1.26 names \emph{-ta} and \emph{-tavat} together).

\subsection*{Your turn}
Say these aloud:

\begin{itemize}
\raggedright
  \item \textquotedblleft{}He came home\textquotedblright{} — \emph{saḥ gṛham āgatavān}
  \item \textquotedblleft{}She came home\textquotedblright{} — \emph{sā gṛham āgatavatī}
  \item \textquotedblleft{}They came home\textquotedblright{} — \emph{te gṛham āgatavantaḥ}
\end{itemize}

\begin{tcolorbox}[breakable,colback=teal!4,colframe=teal!35!black,title={Before you move on}]
\textquotedblleft{}She went\textquotedblright{}? (\textbf{\sk{सा} \sk{गतवती}}.) \textquotedblleft{}I did the work\textquotedblright{}, said by a man? (\textbf{\sk{अहं} \sk{कार्यं} \sk{कृतवान्}}.) Who decides \emph{-vān} or \emph{-vatī}? (The one who did it.)
\end{tcolorbox}
\section[\texorpdfstring{\sk{सः} \sk{पठति} \sk{स्म} (saḥ paṭhati sma)}{सः पठति स्म (saḥ paṭhati sma)}]{\sk{सः} \sk{पठति} \sk{स्म} (saḥ paṭhati sma) — what was going on}

\label{lesson:SA-C276-sah-pathati-sma}



\begin{quote}
Say \textquotedblleft{}he read the book\textquotedblright{}. (\emph{saḥ pustakaṃ paṭhitavān}.) That tells one finished event. This lesson tells the scene around it.
\end{quote}

\subsection*{You'll want to know: the present with \sk{स्म}}
\begin{quote}
\textbf{\sk{सः} \sk{पठति} \sk{स्म।}}
\end{quote}

\begin{quote}
\emph{saḥ paṭhati sma.} — \textquotedblleft{}He was reading.\textquotedblright{}
\end{quote}

\textbf{\sk{स्म}} (\emph{sma}) after a present puts it in the past: \textquotedblleft{}he was reading\textquotedblright{}, or \textquotedblleft{}he used to read\textquotedblright{}. The verb keeps its present ending. Lanman's vocabulary describes it as a present that tells the past (he cites Whitney §778).

\begin{quote}
\textbf{\sk{बालाः} \sk{क्रीडन्ति} \sk{स्म।}}
\end{quote}

\begin{quote}
\emph{bālāḥ krīḍanti sma.} — \textquotedblleft{}The children were playing.\textquotedblright{}
\end{quote}

\begin{quote}
\textbf{\sk{सा} \sk{प्रतिदिनं} \sk{पठति} \sk{स्म।}}
\end{quote}

\begin{quote}
\emph{sā pratidinaṃ paṭhati sma.} — \textquotedblleft{}She used to read every day.\textquotedblright{}
\end{quote}

\subsection*{You'll want to know: \sk{यदा} … \sk{तदा}}
\textbf{\sk{यदा}} (\emph{yadā}) is \textquotedblleft{}when\textquotedblright{}, opening the scene; \textbf{\sk{तदा}} (\emph{tadā}), \textquotedblleft{}then\textquotedblright{}, opens what broke into it. The scene takes \textbf{\sk{स्म}}, the event the past:

\begin{quote}
\textbf{\sk{यदा} \sk{अहं} \sk{पठामि} \sk{स्म},}
\end{quote}

\begin{quote}
\emph{yadā ahaṃ paṭhāmi sma,}
\end{quote}

\begin{quote}
\textbf{\sk{तदा} \sk{मीरा} \sk{आगतवती।}}
\end{quote}

\begin{quote}
\emph{tadā mīrā āgatavatī.} — \textquotedblleft{}When I was reading, Mira came.\textquotedblright{}
\end{quote}

\begin{quote}
\textbf{\sk{यदा} \sk{सः} \sk{खादति} \sk{स्म},}
\end{quote}

\begin{quote}
\emph{yadā saḥ khādati sma,}
\end{quote}

\begin{quote}
\textbf{\sk{तदा} \sk{अहं} \sk{गतवान्।}}
\end{quote}

\begin{quote}
\emph{tadā ahaṃ gatavān.} — \textquotedblleft{}While he was eating, I left.\textquotedblright{}
\end{quote}

\noindent
\begin{tabularx}{\linewidth}{@{}>{\raggedright\arraybackslash}X>{\raggedright\arraybackslash}X@{}}
\toprule
\textbf{the sentence} & \textbf{what it tells} \\
\midrule
\textbf{\sk{सः} \sk{पठितवान्}} & one event: he read \\
\textbf{\sk{सः} \sk{पठति} \sk{स्म}} & the scene: he was reading \\
\bottomrule
\end{tabularx}

\subsection*{Your turn}
Say these aloud:

\begin{itemize}
\raggedright
  \item \textquotedblleft{}I was reading\textquotedblright{} — \emph{ahaṃ paṭhāmi sma}
  \item \textquotedblleft{}I was eating\textquotedblright{} — \emph{ahaṃ khādāmi sma}
  \item \textquotedblleft{}I was writing\textquotedblright{} — \emph{ahaṃ likhāmi sma}
\end{itemize}

\begin{tcolorbox}[breakable,colback=teal!4,colframe=teal!35!black,title={Before you move on}]
\textquotedblleft{}He was playing\textquotedblright{}? (\textbf{\sk{सः} \sk{क्रीडति} \sk{स्म}}.) \textquotedblleft{}When I was reading, Arun came\textquotedblright{}? (\textbf{\sk{यदा} \sk{अहं} \sk{पठामि} \sk{स्म}, \sk{तदा} \sk{अरुणः} \sk{आगतवान्}}.) Which form tells the scene? (The present with \textbf{\sk{स्म}}.)
\end{tcolorbox}
\section[\texorpdfstring{\sk{घटते} (ghaṭate)}{घटते (ghaṭate)}]{\sk{घटते} (ghaṭate) — happens}

\label{lesson:SA-C276-ghatate}



\begin{quote}
Before the new one: say \textquotedblleft{}when I was reading, Mira came\textquotedblright{}. (\emph{yadā ahaṃ paṭhāmi sma, tadā mīrā āgatavatī}.)
\end{quote}

\subsection*{You'll want to know: \sk{घटते}}
\textbf{\sk{घटते}} — \emph{ghaṭate} — \textquotedblleft{}happens\textquotedblright{}. It ends in \emph{-te}, as \textbf{\sk{रोचते}} (\emph{rocate}) does.

\begin{quote}
\textbf{\sk{अत्र} \sk{किं} \sk{घटते}?}
\end{quote}

\begin{quote}
\emph{atra kiṃ ghaṭate?} — \textquotedblleft{}What is happening here?\textquotedblright{}
\end{quote}

\begin{quote}
\textbf{\sk{एतत्} \sk{प्रतिदिनं} \sk{घटते।}}
\end{quote}

\begin{quote}
\emph{etat pratidinaṃ ghaṭate.} — \textquotedblleft{}This happens every day.\textquotedblright{}
\end{quote}

\textquotedblleft{}It happened\textquotedblright{} is its \emph{-ta} form, \textbf{\sk{घटितम्}} (\emph{ghaṭitam}), the way \textbf{\sk{गतः}} says \textquotedblleft{}gone\textquotedblright{}:

\begin{quote}
\textbf{\sk{ह्यः} \sk{किं} \sk{घटितम्}?}
\end{quote}

\begin{quote}
\emph{hyaḥ kiṃ ghaṭitam?} — \textquotedblleft{}What happened yesterday?\textquotedblright{}
\end{quote}

\begin{cousinweb}
\textbf{\sk{घटते}} (\emph{ghaṭate}) is built on the root √ghaṭ. Monier-Williams gives its first sense as being busy with something or striving for it, and a later one as happening, taking place; he gives \emph{ghaṭita} as \textquotedblleft{}happened\textquotedblright{}. It lives on as Marathi \emph{ghaḍṇẽ}, \textquotedblleft{}to happen\textquotedblright{} (CDIAL 4407).
\end{cousinweb}

\subsection*{Your turn}
Say these aloud:

\begin{itemize}
\raggedright
  \item \emph{ghaṭate}
  \item \textquotedblleft{}What is happening here?\textquotedblright{} — \emph{atra kiṃ ghaṭate?}
  \item \textquotedblleft{}What happened yesterday?\textquotedblright{} — \emph{hyaḥ kiṃ ghaṭitam?}
\end{itemize}

\begin{tcolorbox}[breakable,colback=teal!4,colframe=teal!35!black,title={Before you move on}]
What does \textbf{\sk{घटते}} mean? (\textquotedblleft{}Happens\textquotedblright{}.) \textquotedblleft{}What happened?\textquotedblright{} (\textbf{\sk{किं} \sk{घटितम्}?})
\end{tcolorbox}
\section[\texorpdfstring{\sk{घ} (gha)}{घ (gha)}]{\sk{घ} — one character, met first in \sk{घटते}}

\label{lesson:SA-S226-letter-gha}



\begin{quote}
Before the new one: say \textbf{\sk{घटते}} (\emph{ghaṭate}), \textquotedblleft{}happens\textquotedblright{}. Then say the sounds of \textbf{\sk{ग}} and \textbf{\sk{ध}}. (\emph{ga}, \emph{dha}.)

One character this time: the one that opens \textbf{\sk{घटते}}.
\end{quote}

\subsection*{Script you'll notice: \sk{घ}}
\textbf{\sk{घ}} — \emph{gha}.

It is a \textbf{consonant}: \emph{g} with a breath after it, as \textbf{\sk{ध}} is \emph{d} with a breath.

What it is made of:

\begin{itemize}
\raggedright
  \item a continuous upper curl, middle hook, lower bowl, and rising right side
  \item a short descending lower stem
  \item the top shirorekhā
\end{itemize}

You have met it here:

\begin{itemize}
\raggedright
  \item \textbf{\sk{घटते}} \emph{ghaṭate} — happens
\end{itemize}

\subsection*{Writing: \sk{घ}}
\hlblockfigure{\detokenize{figures/SA-S226-letter-gha-filmstrip.pdf}}{How \sk{घ} is written, stroke by stroke}

\begin{itemize}
\raggedright
  \item \textbf{1.} start below the headline at the upper-left, curve down and clockwise around the upper curl, turn back left through the middle hook, continue clockwise around the lower bowl, and rise along the right side to the headline
  \item \textbf{2.} without lifting, come back down the right side to the lower junction
  \item \textbf{3.} descend the short stem below the bowl
  \item \textbf{4.} lift and draw the top shirorekhā left-to-right
\end{itemize}

\textbf{Pen lifts: 1.} The pen comes up 1 time and no more.

\begin{quote}
The source shows three movements; native writers draw \sk{घ} as two pen-down strokes (83\% of HP Labs India's native-writer samples, 68 of 82). This ductus follows the three movements in order and direction, lifting only before the headline: after rising to the headline the pen comes back down the right side into the short lower stem.
\end{quote}

\begin{quote}
This is one attested teaching order and not a national standard — handwriting here is taught with school-to-school variation. Source: Opiaterein, ‘Deva-\sk{घ}-order.gif’, strokes 1–3, Wikimedia Commons, 9 May 2009.
\end{quote}

\subsection*{Your turn}
\begin{itemize}
\raggedright
  \item \emph{Look:} at this, and find \sk{घ} in it
\end{itemize}

\begin{quote}
\sk{घटते}
\end{quote}

\begin{itemize}
\raggedright
  \item \emph{Trace it:} \sk{घ} three times, saying \emph{gha} as you finish each one
  \item \emph{Look:} at \textbf{\sk{किं} \sk{घटितम्}?} and find \sk{घ} once more
\end{itemize}

\begin{tcolorbox}[breakable,colback=teal!4,colframe=teal!35!black,title={Before you move on}]
Which character is this — \sk{घ}? What sound does it carry? (\emph{gha}.) Which word did it come from? (\textbf{\sk{घटते}}, \textquotedblleft{}happens\textquotedblright{}.)
\end{tcolorbox}
\section[(review)]{Review — what happened, and what was going on}

\label{lesson:SA-R276-what-happened}



\begin{quote}
Say \textquotedblleft{}they will go\textquotedblright{} and \textquotedblleft{}I will see the city\textquotedblright{}. (\emph{te gamiṣyanti}, \emph{ahaṃ nagaraṃ drakṣyāmi}.)
\end{quote}

\subsection*{Your turn}
\begin{itemize}
\raggedright
  \item \emph{Say it:} Yesterday Mira went to the shop.
  \item \emph{Say it:} Arun ate the food, and Mira read the book.
  \item \emph{Say it:} The children were playing.
  \item \emph{Say it:} When I was reading, Mira came.
  \item \emph{Say it:} What happened yesterday?
  \item \emph{Read it:} \textbf{\sk{घटते}} and \textbf{\sk{किं} \sk{घटितम्}?} aloud
\end{itemize}

\begin{tcolorbox}[breakable,colback=teal!4,colframe=teal!35!black,title={Before you move on}]
\textbf{\sk{गतवती}}: a man or a woman? (A woman.) \textbf{\sk{सः} \sk{पठति} \sk{स्म}}: one event, or the scene? (The scene.) Now the other way round: \textquotedblleft{}she read the book\textquotedblright{}? (\textbf{\sk{सा} \sk{पुस्तकं} \sk{पठितवती}}.) \textquotedblleft{}What did he say?\textquotedblright{} (\textbf{\sk{सः} \sk{किम्} \sk{उक्तवान्}?})
\end{tcolorbox}
